\chapter{Xây Dựng Dữ Liệu và Phân Tích Khám Phá}
\label{chap:data}

Chương này mô tả toàn bộ pipeline xây dựng dữ liệu mô phỏng và phân tích khám phá (EDA) nhằm hiểu rõ đặc tính của bộ dữ liệu trước khi huấn luyện mô hình Knowledge Tracing.

% ─────────────────────────────────────────────────────────────────────────────
\section{Tổng Quan Quy Trình Notebooks}
\label{sec:notebook_pipeline}

Toàn bộ phần huấn luyện Knowledge Tracing được tổ chức thành \textbf{bốn notebook chạy tuần tự}, mỗi notebook nhận trực tiếp đầu ra của notebook trước (Hình~\ref{fig:notebook_pipeline}). Quy trình này tách rõ bốn giai đoạn: khám phá dữ liệu, kỹ thuật đặc trưng, huấn luyện và đánh giá; nhờ đó có thể tái lập từng bước một cách độc lập. Toàn bộ pipeline cố định \texttt{seed = 42}; hai notebook huấn luyện và đánh giá (\texttt{03}, \texttt{04}) chạy trên Kaggle với GPU NVIDIA Tesla T4 và PyTorch 2.6.

\begin{figure}[H]
  \centering
  \begin{tikzpicture}[
    node distance=0.7cm and 0.9cm,
    nb/.style={rectangle, rounded corners, draw, thick, fill=blue!8,
      text width=3.1cm, minimum height=1.25cm, align=center, font=\small},
    io/.style={rectangle, draw, fill=gray!10, text width=3.0cm,
      minimum height=0.9cm, align=center, font=\scriptsize\ttfamily},
    arr/.style={-{Stealth[length=2.2mm]}, thick}
  ]
    \node[io] (raw) {data/raw/*};
    \node[nb, below=of raw] (n1) {\textbf{01\_eda}\\Phân tích khám phá, phát hiện \& fix issue};
    \node[nb, below=of n1] (n2) {\textbf{02\_feature\_engineering}\\Tạo nhãn, encode, build sequences};
    \node[nb, below=of n2] (n3) {\textbf{03\_train\_models}\\Ablation 8 cấu hình KT};
    \node[nb, below=of n3] (n4) {\textbf{04\_evaluation}\\Phân tích sâu, chọn mô hình};

    \node[io, right=of n2] (proc) {data/processed/ *.pkl};
    \node[io, right=of n3] (mdl) {results/models/ metrics/ ablation/};
    \node[io, right=of n4] (plt) {results/plots/*};

    \draw[arr] (raw) -- (n1);
    \draw[arr] (n1) -- (n2);
    \draw[arr] (n2) -- (n3);
    \draw[arr] (n3) -- (n4);
    \draw[arr] (n2) -- (proc);
    \draw[arr] (n3) -- (mdl);
    \draw[arr] (n4) -- (plt);
  \end{tikzpicture}
  \caption{Pipeline bốn notebook chạy tuần tự, từ dữ liệu thô đến mô hình KT triển khai.}
  \label{fig:notebook_pipeline}
\end{figure}

Vai trò tóm tắt của từng notebook như sau:

\begin{itemize}
  \item \textbf{\texttt{01\_eda}} --- Kiểm tra chất lượng dữ liệu thô (null, mất cân bằng, bất thường) và phát hiện các vấn đề cần sửa \emph{trước khi} feature engineering: chuẩn hóa quality thành nhãn nhị phân (cho AUC) và nhãn float (cho RMSE), cảnh báo độ dài chuỗi đồng đều (uniform), và đối chiếu tự đánh giá với năng lực thực để phát hiện confidence bias ($r \approx 0.72$). Chi tiết trình bày ở các Mục~\ref{sec:eda_qbank}--\ref{sec:eda_bias}.
  \item \textbf{\texttt{02\_feature\_engineering}} --- Biến dữ liệu đã làm sạch thành các chuỗi sẵn sàng huấn luyện: tạo nhãn \texttt{correct} $\in\{0,1\}$ và \texttt{quality\_norm} $\in\{0,0.5,1\}$, encode KC/user/question/difficulty, chia tập \emph{theo người dùng} tỷ lệ 70/15/15 để tránh data leakage, và build sequence cho mỗi user. Đầu ra là \texttt{train/val/test\_seqs.pkl}.
  \item \textbf{\texttt{03\_train\_models}} --- Huấn luyện và so sánh tám cấu hình KT (Baseline, BKT, DKT/SAKT theo nhãn binary và quality, cùng biến thể deep) trên cùng dữ liệu và thiết lập. Chi tiết ablation ở Mục~\ref{sec:configs} và Chương~\ref{chap:experiment}.
  \item \textbf{\texttt{04\_evaluation}} --- Phân tích sâu sau huấn luyện (per-KC AUC, theo vị trí chuỗi, theo độ khó, theo loại người dùng, cold-start, phân tích lỗi) để rút kết luận có phản biện và chọn DKT-Quality$\star$ làm mô hình triển khai (Mục~\ref{sec:deep_analysis}).
\end{itemize}

Quy trình chi tiết của từng giai đoạn được trình bày trong phần còn lại của chương này (notebook \texttt{01}--\texttt{02}) và Chương~\ref{chap:experiment} (notebook \texttt{03}--\texttt{04}); các bước tái lập đầy đủ được liệt kê trong Phụ Lục~\ref{apx:reproducibility}.

\medskip
\noindent\textit{Lưu ý số liệu:} các con số trong báo cáo dùng snapshot cuối (1.578 câu hỏi, 1.000 người dùng, $\approx$40.140 tương tác). Một số bản nháp notebook đầu còn ghi snapshot cũ hơn (1.175 câu / 500 user / 20.000 tương tác); pipeline và các bước xử lý là như nhau.

% ─────────────────────────────────────────────────────────────────────────────
\section{Thu Thập và Kiểm Định Question Bank}
\label{sec:qbank_build}

Khác với phần lớn nghiên cứu KT chỉ dùng dataset có sẵn, khóa luận tự xây dựng question bank phỏng vấn IT qua một quy trình lai (hybrid) gồm bốn giai đoạn: thu thập, sinh tự động, kiểm định chất lượng, và chuẩn hóa.

\subsection{Thu thập từ nguồn công khai}
Câu hỏi được thu thập (scraping) từ nhiều kho câu hỏi phỏng vấn mở, sau đó trích dẫn nguồn: bộ câu hỏi Data Engineer (Obenner), bộ câu hỏi SQL (LearningZone), và các bộ câu hỏi DS/DA cộng đồng (Youssef, Alexey). Mỗi câu được ánh xạ sơ bộ về vai trò (DA/DS/DE) và nhóm kỹ năng (\texttt{skill\_groups}).

\subsection{Sinh câu hỏi bằng LLM}
Để tăng độ phủ các KC còn thiếu, một pipeline đa mô hình được sử dụng: \textbf{Gemini Flash} sinh câu hỏi mới theo từng KC và mức độ; mô hình mạnh hơn tạo \textit{reference answer} và \textit{evaluation points}; một mô hình \textit{judge} chấm chất lượng và một mô hình \textit{enhancer} tinh chỉnh diễn đạt. Bốn loại câu hỏi khách quan (MC\_SINGLE, TRUE\_FALSE, FILL\_BLANK) và hai loại tự luận (THEORY, CODING) được sinh kèm metadata đầy đủ.

\subsection{Kiểm định chất lượng (QC)}
Mỗi câu hỏi đi qua: (1) \textbf{kiểm tra schema} theo \texttt{question\_bank.schema.json} (bắt buộc các trường \texttt{question\_id}, \texttt{question\_text}, \texttt{roles}, \texttt{difficulty\_label}, \texttt{skill\_groups}, \texttt{answers}); (2) \textbf{QC bằng RAG} đối chiếu nội dung với tài liệu tham chiếu để phát hiện câu sai sự thật; (3) \textbf{human review} trên mẫu ngẫu nhiên; và (4) \textbf{chuẩn hóa skill\_groups} -- ánh xạ các tag chi tiết (ví dụ \texttt{SQL\_WINDOW\_FUNCTION}) về 17 KC chuẩn qua bảng tra \texttt{SKILL\_GROUP\_TO\_KC} (82 ánh xạ).

\subsection{Kết quả}
Snapshot dùng cho thực nghiệm KT gồm \textbf{1.578 câu} (17 KC, 3 mức độ). Sau khi tiếp tục thu thập và sinh bổ sung, question bank của ứng dụng triển khai mở rộng lên \textbf{2.122 câu}. Quá trình QC cũng phát hiện và xử lý một số lỗi dữ liệu điển hình (ví dụ câu bị tách rời tiêu đề khỏi nội dung khi scraping), minh họa tầm quan trọng của bước kiểm định.

\section{Pipeline Xây Dựng Dữ Liệu}
\label{sec:data_pipeline}

Do thiếu dữ liệu thực từ người dùng thật, luận văn sử dụng pipeline mô phỏng ba bước để tạo bộ dữ liệu training:

\begin{enumerate}
  \item \textbf{Xây dựng question bank}: Thu thập và QC câu hỏi phỏng vấn qua pipeline 3 mô hình LLM (Gemini Flash sinh câu hỏi, GPT-4o làm reference answer, Claude Haiku judge chất lượng, Claude Sonnet enhance). Kết quả của snapshot dùng cho thực nghiệm KT: 1{,}578 câu hỏi unique, phân loại theo 3 role (DA/DS/DE), 3 mức độ (EASY/MEDIUM/HARD) và 17 Knowledge Components (KC). Về sau, question bank được mở rộng lên 2{,}122 câu cho ứng dụng triển khai (xem Chương~\ref{chap:system}).

  \item \textbf{Tạo virtual users}: Sinh 1{,}000 người dùng ảo với skill vector theo phân phối correlated Gaussian, phản ánh tương quan thực tế giữa các competency (ví dụ: DA có $r(\text{SQL}, \text{Analytics}) \approx 0.55$; DS có $r(\text{Algorithm}, \text{Evaluation}) \approx 0.70$).

  \item \textbf{Mô phỏng interaction}: Với mỗi người dùng, sinh chuỗi tương tác dựa trên:
  \begin{itemize}
    \item \textbf{IRT}: $P(\text{correct}) = \sigma\bigl((\theta - b) \times 6\bigr)$, trong đó $\theta$ là skill hiện tại và $b$ là độ khó câu hỏi.
    \item \textbf{Power-based graded response}: $P(Q=2) = p^{\alpha}$, $P(Q=0) = (1-p)^{\alpha}$, $\alpha = 2.0$.
    \item \textbf{Skill update}: $\theta \mathrel{+}= \eta(1 - \theta)$ khi đúng; $\theta \mathrel{-}= \gamma\theta$ khi sai, với $\eta = 0.08$, $\gamma = 0.002$.
  \end{itemize}
\end{enumerate}

% ─────────────────────────────────────────────────────────────────────────────
\section{Phân Tích Question Bank}
\label{sec:eda_qbank}

\begin{figure}[H]
  \centering
  \includegraphics[width=\textwidth]{eda_qbank_distribution.png}
  \caption{Phân bố câu hỏi theo role và mức độ khó trong question bank}
  \label{fig:qbank_dist}
\end{figure}

Hình~\ref{fig:qbank_dist} cho thấy question bank có phân bố khá đồng đều giữa ba role (DA/DS/DE), với mỗi role có khoảng 500--550 câu hỏi. Trong mỗi role, câu hỏi MEDIUM chiếm tỷ lệ lớn nhất, phù hợp với thực tế phỏng vấn (câu hỏi quá dễ hoặc quá khó đều ít giá trị phân biệt).

\begin{figure}[H]
  \centering
  \includegraphics[width=\textwidth]{eda_qbank_skill_coverage.png}
  \caption{Độ phủ câu hỏi theo từng Knowledge Component (KC)}
  \label{fig:qbank_skill}
\end{figure}

Hình~\ref{fig:qbank_skill} thể hiện số lượng câu hỏi theo từng KC. Có sự chênh lệch rõ rệt: một số KC như \texttt{SQL\_DATABASE} (DA) và \texttt{ALGORITHM\_THEORY} (DS) có số câu hỏi nhiều hơn, phản ánh tầm quan trọng và độ phổ biến của các chủ đề này trong thực tế phỏng vấn IT.

% ─────────────────────────────────────────────────────────────────────────────
\section{Phân Tích Người Dùng}
\label{sec:eda_users}

\begin{figure}[H]
  \centering
  \includegraphics[width=\textwidth]{eda_user_distribution.png}
  \caption{Phân bố người dùng theo role (trái) và trình độ (phải)}
  \label{fig:user_dist}
\end{figure}

Hình~\ref{fig:user_dist} (trái) xác nhận phân bổ đồng đều theo role: DA = 334, DS = 333, DE = 333. Biểu đồ phải cho thấy phân phối trình độ thực tế sau khi phân loại dựa trên skill vector: intermediate = 532 (53.2\%), beginner = 409 (40.9\%), advanced = 59 (5.9\%). Phân phối này lệch so với mục tiêu thiết kế (35\%/50\%/15\%) do quá trình phân loại lại (\textit{reclassification}) dựa trên skill vector thực tế sau sampling -- nhiều người dùng thuộc tier advanced bị reclassify về intermediate khi mean skill nằm trong khoảng $[0.40, 0.70]$.

\begin{figure}[H]
  \centering
  \includegraphics[width=\textwidth]{eda_user_skill_vector.png}
  \caption{Heatmap skill vector của 50 người dùng mẫu (mỗi hàng là 1 user)}
  \label{fig:user_skill}
\end{figure}

Hình~\ref{fig:user_skill} minh họa sự đa dạng của skill vector giữa các người dùng. Mỗi hàng đại diện cho một người dùng; màu sắc phản ánh mức skill từ thấp (xanh lam) đến cao (vàng/đỏ). Có thể quan sát tương quan dương giữa các KC trong cùng role -- người dùng có skill cao ở một KC thường có skill trung bình trở lên ở các KC liên quan.

% ─────────────────────────────────────────────────────────────────────────────
\section{Phân Tích Tương Tác}
\label{sec:eda_interaction}

\subsection{Phân phối quality score}

\begin{figure}[H]
  \centering
  \includegraphics[width=0.8\textwidth]{eda_quality_distribution.png}
  \caption{Phân phối quality score (0 = Fail, 1 = Pass HR, 2 = Pass Tech)}
  \label{fig:quality_dist}
\end{figure}

Hình~\ref{fig:quality_dist} cho thấy phân phối quality score với Q=2 chiếm tỷ lệ cao nhất ($\sim$62\%), Q=1 khoảng 20\%, và Q=0 ($\sim$18\%). Phân phối này phản ánh mô hình power-based graded response với $\alpha = 2.0$: bucket Q=1 (partial credit) rộng nhất ở vùng skill trung bình, trong khi Q=2 và Q=0 tập trung ở hai đầu phổ skill.

\subsection{Độ dài chuỗi tương tác}

\begin{figure}[H]
  \centering
  \includegraphics[width=0.8\textwidth]{eda_sequence_length.png}
  \caption{Phân phối độ dài chuỗi tương tác mỗi người dùng}
  \label{fig:seq_len}
\end{figure}

Hình~\ref{fig:seq_len} cho thấy toàn bộ người dùng có đúng 40 tương tác. Mặc dù pipeline thiết kế randomize 20--60 câu/user, trên thực tế mỗi role chỉ có $\sim$40 câu trong pool sau bước stratified sampling theo competency, dẫn đến sequence length bị cap tại pool size. Đây là hạn chế của quy mô question bank hiện tại (xem Mục~\ref{sec:limitations} để thảo luận chi tiết).

\subsection{Quality theo mức độ khó}

\begin{figure}[H]
  \centering
  \includegraphics[width=0.85\textwidth]{eda_quality_by_difficulty.png}
  \caption{Phân phối quality score theo độ khó câu hỏi (EASY/MEDIUM/HARD)}
  \label{fig:quality_diff}
\end{figure}

Hình~\ref{fig:quality_diff} xác nhận tính nhất quán của mô hình IRT: câu hỏi EASY có tỷ lệ Q=2 cao hơn rõ rệt, câu HARD có tỷ lệ Q=0 cao hơn. Câu hỏi MEDIUM có phân phối cân bằng nhất giữa ba mức -- đây là vùng có độ phân biệt (discrimination) cao nhất trong IRT, phù hợp nhất cho training.

\subsection{Quality heatmap theo competency}

\begin{figure}[H]
  \centering
  \includegraphics[width=\textwidth]{eda_quality_heatmap_competency.png}
  \caption{Tỷ lệ Pass (Q$\geq$1) trung bình theo competency và role}
  \label{fig:quality_heatmap}
\end{figure}

Hình~\ref{fig:quality_heatmap} cho thấy tỷ lệ pass khác nhau đáng kể giữa các competency. Trong nhóm DA, \texttt{SQL\_DATABASE} có pass rate cao hơn \texttt{STATISTICS\_EXPERIMENTATION} -- phản ánh penalty mà pipeline thiết kế (người DA thường yếu statistics hơn SQL). Các KC của DS như \texttt{DEEP\_LEARNING} và \texttt{NLP} có pass rate thấp nhất, nhất quán với độ phức tạp cao của các chủ đề này.

% ─────────────────────────────────────────────────────────────────────────────
\section{Phân Tích Sparsity và Learning Curve}
\label{sec:eda_sparsity}

\subsection{Ma trận sparsity user--KC}

\begin{figure}[H]
  \centering
  \includegraphics[width=\textwidth]{eda_sparsity_matrix.png}
  \caption{Ma trận tương tác user--KC: màu tối = có tương tác, trắng = không có}
  \label{fig:sparsity}
\end{figure}

Hình~\ref{fig:sparsity} minh họa ma trận user--KC. Do mỗi user chỉ trả lời câu hỏi thuộc role của mình, matrix có cấu trúc block diagonal rõ ràng: DA users chỉ có tương tác với KC của DA, tương tự cho DS và DE. Điều này dẫn đến sparsity cao khi xét toàn bộ không gian KC ($\sim$77\% entries rỗng), nhưng sparsity trong từng role thấp hơn đáng kể ($\sim$15--20\%), đủ để BKT hoạt động ổn định.

\subsection{Learning curve tổng thể}

\begin{figure}[H]
  \centering
  \includegraphics[width=0.85\textwidth]{eda_learning_curve.png}
  \caption{Learning curve: skill trung bình theo vị trí trong chuỗi tương tác}
  \label{fig:lc_overall}
\end{figure}

Hình~\ref{fig:lc_overall} xác nhận rằng pipeline mô phỏng có learning dynamics đúng: skill trung bình tăng dần theo thứ tự câu hỏi trong session. Xu hướng tăng phản ánh tham số học $\eta = 0.08$ và forgetting rate thấp $\gamma = 0.002$, tạo ra learning curve có hình dạng đặc trưng của mô hình exponential growth with saturation.

\subsection{Learning curve theo loại người dùng}

\begin{figure}[H]
  \centering
  \includegraphics[width=\textwidth]{eda_learning_curve_by_type.png}
  \caption{Learning curve phân tầng theo trình độ người dùng (beginner/intermediate/advanced)}
  \label{fig:lc_bytype}
\end{figure}

Hình~\ref{fig:lc_bytype} cho thấy ba nhóm có điểm xuất phát và tốc độ học khác nhau rõ rệt. Người dùng \textit{advanced} bắt đầu ở mức skill cao và có ít dư địa cải thiện (ceiling effect). Ngược lại, nhóm \textit{beginner} có tốc độ tăng skill tuyệt đối lớn nhất -- phù hợp với đặc tính của công thức $\Delta\theta = \eta(1 - \theta)$: skill tăng nhanh hơn khi còn thấp.

% ─────────────────────────────────────────────────────────────────────────────
\section{Phân Tích Confidence Bias}
\label{sec:eda_bias}

\begin{figure}[H]
  \centering
  \includegraphics[width=0.75\textwidth]{eda_confidence_bias.png}
  \caption{Self-rating so với skill thực: phát hiện confidence bias}
  \label{fig:confidence_bias}
\end{figure}

Hình~\ref{fig:confidence_bias} phân tích mối quan hệ giữa self-rating (điểm tự đánh giá) và skill thực tế (được mô phỏng). Hệ số tương quan $r \approx 0.72$ cho thấy self-rating có giá trị thông tin nhất định nhưng không hoàn hảo. Nhóm beginner có xu hướng \textit{overconfidence} (self-rating cao hơn skill thực), trong khi nhóm advanced thỉnh thoảng có \textit{underconfidence} -- phản ánh các bias tâm lý được ghi nhận trong tài liệu tự đánh giá năng lực.

% ─────────────────────────────────────────────────────────────────────────────
\section{Tóm Tắt Đặc Tính Dữ Liệu}
\label{sec:data_summary}

Bảng~\ref{tab:eda_summary} tổng hợp các đặc tính chính của bộ dữ liệu mô phỏng:

\begin{table}[H]
  \centering
  \caption{Tóm tắt đặc tính bộ dữ liệu EDA}
  \label{tab:eda_summary}
  \begin{tabular}{ll}
    \toprule
    \textbf{Đặc tính} & \textbf{Giá trị / Nhận xét} \\
    \midrule
    Tổng câu hỏi & 1{,}578 (DA: $\sim$525, DS: $\sim$544, DE: $\sim$509) \\
    Tổng người dùng & 1{,}000 (DA: 334, DS: 333, DE: 333) \\
    Trình độ (thực tế) & Beginner 40.9\%, Intermediate 53.2\%, Advanced 5.9\% \\
    Tổng tương tác & $\sim$40{,}000 (trung bình 40/user) \\
    Sequence length & Hằng định = 40 (do pool cap) \\
    Quality distribution & Q=2: $\sim$62\%, Q=1: $\sim$20\%, Q=0: $\sim$18\% \\
    Pass rate (Q$\geq$1) & $\sim$82\% \\
    Sparsity (user--KC) & $\sim$77\% toàn bộ; $\sim$15--20\% trong từng role \\
    Learning dynamics & Skill tăng đơn điệu theo session order ($\Delta > 0$) \\
    Confidence bias & $r(\text{self-rating}, \text{skill}) \approx 0.72$ \\
    \bottomrule
  \end{tabular}
\end{table}

Bộ dữ liệu có learning dynamics hợp lý và phân phối quality nhất quán với thiết kế IRT. Hai hạn chế chính -- sequence length hằng định và phân phối advanced thấp -- được thảo luận chi tiết trong phần Hạn chế (Mục~\ref{sec:limitations} của Chương~\ref{chap:experiment}).
